\clearpage
\SourceLeaf{149}
\begin{center}
[IV.]

{\large [CRITICAL REMARKS ON THE METHOD OF LEAST SQUARES.]}
\end{center}



\Needspace{6\baselineskip}
\begin{center}
\textsc{Gauss} to \textsc{Olbers}. Göttingen, 22 February 1819.
\end{center}

\noindent . . . . . . . . . . . I am also now working on a new justification of the so-called method of least squares. My first justification assumes that the probability of the observational error $x$ is represented by $e^{-hhxx}$. The method then gives the most probable results with full rigor and in every case. If the law of errors is unknown, it is \textls{impossible} to specify the most probable results from observations \textls{already made}. \textsc{Laplace} approached the matter from a different side and chose a principle that also leads to the\,method\,of\,least\ squares\HistoricalNote{*)}{*)}{Independently of the law of errors, in fact.} when the number of observations is infinitely large. With a moderate number of observations, however, one remains entirely in the dark if the law of errors is unknown. Even \textsc{Laplace} has nothing better to say in this case than that one should apply the method\ of least\ squares here too, because it makes for convenient calculation. I have now found that if one chooses a principle somewhat different from \textsc{Laplace}'s (one that \textls{no one} can deny we have at least as much right to adopt as the one chosen by \textsc{Laplace}, and that, in my opinion, \textls{anyone without a preconceived view} must declare more natural than \textsc{Laplace}'s), one enjoys all those\space%
